\haddockmoduleheading{Control.Monad}
\label{module:Control.Monad}
\haddockbeginheader
{\haddockverb\begin{verbatim}
module Control.Monad (
    Functor(fmap), Monad((>>=), (>>), return), MonadFail(fail),
    MonadPlus(mzero, mplus), mapM, mapM_, forM, forM_, sequence, sequence_,
    (=<<), (>=>), (<=<), forever, void, join, msum, filterM, mapAndUnzipM,
    zipWithM, zipWithM_, foldM, foldM_, replicateM, replicateM_, guard, when,
    unless, liftM, liftM2, liftM3, liftM4, liftM5, ap
  ) where\end{verbatim}}
\haddockendheader

The Control.Monad module provides the \haddockid{Functor}, \haddockid{Monad}, \haddockid{MonadFail} and
 \haddockid{MonadPlus} classes, together with some useful operations on monads.\par
\section{Functor and monad classes}
\begin{haddockdesc}
\item[\begin{tabular}{@{}l}
class Functor f where
\end{tabular}]
{\haddockbegindoc
A type \haddocktt{f} is a Functor if it provides a function \haddocktt{fmap} which, given any types \haddocktt{a} and \haddocktt{b}
lets you apply any function from \haddocktt{(a -> b)} to turn an \haddocktt{f a} into an \haddocktt{f b}, preserving the
structure of \haddocktt{f}. Furthermore \haddocktt{f} needs to adhere to the following:\par
\vbox{\begin{description}
\item[Identity]\hfill \par
\haddocktt{\haddockid{fmap} \haddockid{id} == \haddockid{id}}
\item[Composition]\hfill \par
\haddocktt{\haddockid{fmap} (f . g) == \haddockid{fmap} f . \haddockid{fmap} g}
\end{description}}
Note, that the second law follows from the free theorem of the type \haddockid{fmap} and
the first law, so you need only check that the former condition holds.
See these articles by \href{https://www.schoolofhaskell.com/user/edwardk/snippets/fmap}{School of Haskell} or
\href{https://github.com/quchen/articles/blob/master/second_functor_law.md}{David Luposchainsky}
for an explanation.\par
\haddockpremethods{}\emph{Methods}
\begin{haddockdesc}
\item[\begin{tabular}{@{}l}
fmap :: (a -> b) -> f a -> f b
\end{tabular}]
{\haddockbegindoc
\haddockid{fmap} is used to apply a function of type \haddocktt{(a -> b)} to a value of type \haddocktt{f a},
 where f is a functor, to produce a value of type \haddocktt{f b}.
 Note that for any type constructor with more than one parameter (e.g., \haddocktt{Either}),
 only the last type parameter can be modified with \haddockid{fmap} (e.g., \haddocktt{b} in `Either a b`).\par
Some type constructors with two parameters or more have a \haddocktt{\haddockid{Bifunctor}} instance that allows
 both the last and the penultimate parameters to be mapped over.\par
\subsubsection*{\textbf{Examples}}
Convert from a \haddocktt{\haddockid{Maybe} Int} to a \haddocktt{Maybe String}
 using \haddockid{show}:\par
\begin{quote}
{\haddockverb\begin{verbatim}
>>> fmap show Nothing
Nothing

>>> fmap show (Just 3)
Just "3"

\end{verbatim}}
\end{quote}
Convert from an \haddocktt{\haddockid{Either} Int Int} to an
 \haddocktt{Either Int String} using \haddockid{show}:\par
\begin{quote}
{\haddockverb\begin{verbatim}
>>> fmap show (Left 17)
Left 17

>>> fmap show (Right 17)
Right "17"

\end{verbatim}}
\end{quote}
Double each element of a list:\par
\begin{quote}
{\haddockverb\begin{verbatim}
>>> fmap (*2) [1,2,3]
[2,4,6]

\end{verbatim}}
\end{quote}
Apply \haddockid{even} to the second element of a pair:\par
\begin{quote}
{\haddockverb\begin{verbatim}
>>> fmap even (2,2)
(2,True)

\end{verbatim}}
\end{quote}
It may seem surprising that the function is only applied to the last element of the tuple
 compared to the list example above which applies it to every element in the list.
 To understand, remember that tuples are type constructors with multiple type parameters:
 a tuple of 3 elements \haddocktt{(a,b,c)} can also be written \haddocktt{(,,) a b c} and its \haddocktt{Functor} instance
 is defined for \haddocktt{Functor ((,,) a b)} (i.e., only the third parameter is free to be mapped over
 with \haddocktt{fmap}).\par
It explains why \haddocktt{fmap} can be used with tuples containing values of different types as in the
 following example:\par
\begin{quote}
{\haddockverb\begin{verbatim}
>>> fmap even ("hello", 1.0, 4)
("hello",1.0,True)

\end{verbatim}}
\end{quote}}
\end{haddockdesc}}
\end{haddockdesc}
\begin{haddockdesc}
\item[\begin{tabular}{@{}l}
instance Functor Complex\\instance Functor IO\\instance Functor Maybe\\instance Functor {[}{]}\\instance Functor Array i\\instance Functor -> r\\\end{tabular}]
\end{haddockdesc}
\begin{haddockdesc}
\item[\begin{tabular}{@{}l}
class Applicative m => Monad m where
\end{tabular}]
{\haddockbegindoc
The \haddockid{Monad} class defines the basic operations over a \emph{monad},
a concept from a branch of mathematics known as \emph{category theory}.
From the perspective of a Haskell programmer, however, it is best to
think of a monad as an \emph{abstract datatype} of actions.
Haskell's \haddocktt{do} expressions provide a convenient syntax for writing
monadic expressions.\par
Instances of \haddockid{Monad} should satisfy the following:\par
\vbox{\begin{description}
\item[Left identity]\hfill \par
\haddocktt{\haddockid{return} a \haddockid{>>=} k\ \ =\ \ k a}
\item[Right identity]\hfill \par
\haddocktt{m \haddockid{>>=} \haddockid{return}\ \ =\ \ m}
\item[Associativity]\hfill \par
\haddocktt{m \haddockid{>>=} ({\char '134}x -> k x \haddockid{>>=} h)\ \ =\ \ (m \haddockid{>>=} k) \haddockid{>>=} h}
\end{description}}
Furthermore, the \haddockid{Monad} and \haddockid{Applicative} operations should relate as follows:\par
\vbox{\begin{itemize}
\item
\begin{quote}
{\haddockverb\begin{verbatim}
pure = return\end{verbatim}}
\end{quote}
\item
\begin{quote}
{\haddockverb\begin{verbatim}
m1 <*> m2 = m1 >>= (\x1 -> m2 >>= (\x2 -> return (x1 x2)))\end{verbatim}}
\end{quote}
\end{itemize}}
The above laws imply:\par
\vbox{\begin{itemize}
\item
\begin{quote}
{\haddockverb\begin{verbatim}
fmap f xs  =  xs >>= return . f\end{verbatim}}
\end{quote}
\item
\begin{quote}
{\haddockverb\begin{verbatim}
(>>) = (*>)\end{verbatim}}
\end{quote}
\end{itemize}}
and that \haddockid{pure} and (\haddockid{<*>}) satisfy the applicative functor laws.\par
The instances of \haddockid{Monad} for \haddockid{List}, \haddockid{Maybe} and \haddockid{IO}
defined in the \haddocktt{Prelude} satisfy these laws.\par
\haddockpremethods{}\emph{Methods}
\begin{haddockdesc}
\item[\begin{tabular}{@{}l}
(>>=) :: m a -> (a -> m b) -> m b
\end{tabular}]
{\haddockbegindoc
Sequentially compose two actions, passing any value produced
 by the first as an argument to the second.\par
'\haddocktt{as \haddockid{>>=} bs}' can be understood as the \haddocktt{do} expression\par
\begin{quote}
{\haddockverb\begin{verbatim}
do a <- as
   bs a
\end{verbatim}}
\end{quote}
An alternative name for this function is 'bind', but some people
 may refer to it as 'flatMap', which results from it being equivialent
 to\par
\begin{quote}
{\haddockverb\begin{verbatim}
\x f -> join (fmap f x) :: Monad m => m a -> (a -> m b) -> m b\end{verbatim}}
\end{quote}
which can be seen as mapping a value with
 \haddocktt{Monad m => m a -> m (m b)} and then 'flattening' \haddocktt{m (m b)} to \haddocktt{m b} using \haddockid{join}.\par}
\end{haddockdesc}
\begin{haddockdesc}
\item[\begin{tabular}{@{}l}
(>>) :: m a -> m b -> m b
\end{tabular}]
{\haddockbegindoc
Sequentially compose two actions, discarding any value produced
 by the first, like sequencing operators (such as the semicolon)
 in imperative languages.\par
'\haddocktt{as \haddockid{>>} bs}' can be understood as the \haddocktt{do} expression\par
\begin{quote}
{\haddockverb\begin{verbatim}
do as
   bs
\end{verbatim}}
\end{quote}
or in terms of \haddocktt{\haddockid{(>>=)}} as\par
\begin{quote}
{\haddockverb\begin{verbatim}
as >>= const bs\end{verbatim}}
\end{quote}}
\end{haddockdesc}
\begin{haddockdesc}
\item[\begin{tabular}{@{}l}
return :: a -> m a
\end{tabular}]
{\haddockbegindoc
Inject a value into the monadic type.
 This function should \emph{not} be different from its default implementation
 as \haddockid{pure}. The justification for the existence of this function is
 merely historic.\par}
\end{haddockdesc}}
\end{haddockdesc}
\begin{haddockdesc}
\item[\begin{tabular}{@{}l}
instance Monad Complex\\instance Monad IO\\instance Monad Maybe\\instance Monad {[}{]}\\instance Monad -> r
\end{tabular}]
\end{haddockdesc}
\begin{haddockdesc}
\item[\begin{tabular}{@{}l}
class Monad m => MonadFail m where
\end{tabular}]
{\haddockbegindoc
When a value is bound in \haddocktt{do}-notation, the pattern on the left
 hand side of \haddocktt{<-} might not match. In this case, this class
 provides a function to recover.\par
A \haddockid{Monad} without a \haddockid{MonadFail} instance may only be used in conjunction
 with pattern that always match, such as newtypes, tuples, data types with
 only a single data constructor, and irrefutable patterns (\haddocktt{{\char '176}pat}).\par
Instances of \haddockid{MonadFail} should satisfy the following law: \haddocktt{fail s} should
 be a left zero for \haddockid{>>=},\par
\begin{quote}
{\haddockverb\begin{verbatim}
fail s >>= f  =  fail s
\end{verbatim}}
\end{quote}
If your \haddockid{Monad} is also \haddockid{MonadPlus}, a popular definition is\par
\begin{quote}
{\haddockverb\begin{verbatim}
fail _ = mzero
\end{verbatim}}
\end{quote}
\haddocktt{fail s} should be an action that runs in the monad itself, not an
 exception (except in instances of \haddocktt{MonadIO}).  In particular,
 \haddocktt{fail} should not be implemented in terms of \haddocktt{error}.\par
\haddockpremethods{}\emph{Methods}
\begin{haddockdesc}
\item[\begin{tabular}{@{}l}
fail :: String -> m a
\end{tabular}]
\end{haddockdesc}}
\end{haddockdesc}
\begin{haddockdesc}
\item[\begin{tabular}{@{}l}
instance MonadFail IO\\instance MonadFail Maybe\\instance MonadFail {[}{]}
\end{tabular}]
\end{haddockdesc}
\begin{haddockdesc}
\item[\begin{tabular}{@{}l}
class (Alternative m, Monad m) => MonadPlus m where
\end{tabular}]
{\haddockbegindoc
Monads that also support choice and failure.\par
\haddockpremethods{}\emph{Methods}
\begin{haddockdesc}
\item[\begin{tabular}{@{}l}
mzero :: m a
\end{tabular}]
{\haddockbegindoc
The identity of \haddockid{mplus}.  It should also satisfy the equations\par
\begin{quote}
{\haddockverb\begin{verbatim}
mzero >>= f  =  mzero
v >> mzero   =  mzero\end{verbatim}}
\end{quote}
The default definition is\par
\begin{quote}
{\haddockverb\begin{verbatim}
mzero = empty
\end{verbatim}}
\end{quote}}
\end{haddockdesc}
\begin{haddockdesc}
\item[\begin{tabular}{@{}l}
mplus :: m a -> m a -> m a
\end{tabular}]
{\haddockbegindoc
An associative operation. The default definition is\par
\begin{quote}
{\haddockverb\begin{verbatim}
mplus = (<|>)
\end{verbatim}}
\end{quote}}
\end{haddockdesc}}
\end{haddockdesc}
\begin{haddockdesc}
\item[\begin{tabular}{@{}l}
\end{tabular}]
\end{haddockdesc}
\begin{haddockdesc}
\item[\begin{tabular}{@{}l}
instance MonadPlus IO
\end{tabular}]
{\haddockbegindoc
Takes the first non-throwing \haddockid{IO} action's result.
 \haddockid{mzero} throws an exception.\par}
\end{haddockdesc}
\begin{haddockdesc}
\item[\begin{tabular}{@{}l}
instance MonadPlus Maybe
\end{tabular}]
{\haddockbegindoc
Picks the leftmost \haddockid{Just} value, or, alternatively, \haddockid{Nothing}.\par}
\end{haddockdesc}
\begin{haddockdesc}
\item[\begin{tabular}{@{}l}
instance MonadPlus {[}{]}
\end{tabular}]
{\haddockbegindoc
Combines lists by concatenation, starting from the empty list.\par}
\end{haddockdesc}
\begin{haddockdesc}
\item[\begin{tabular}{@{}l}
\end{tabular}]
\end{haddockdesc}
\section{Functions}
\subsection{Naming conventions}
The functions in this library use the following naming conventions:\par
\vbox{\begin{itemize}
\item
A postfix ’M’ always stands for a function in the Kleisli category:
     The monad type constructor m is added to function results (modulo
     currying) and nowhere else. So, for example,\par
\begin{quote}
{\haddockverb\begin{verbatim}
filter :: (a -> Bool) -> [a] -> [a]
filterM :: (Monad m) => (a -> m Bool) -> [a] -> m [a]
\end{verbatim}}
\end{quote}
\item
A postfix ’{\char '137}’ changes the result type from (m a) to (m ()). Thus,
     for example:\par
\begin{quote}
{\haddockverb\begin{verbatim}
sequence :: Monad m => [m a] -> m [a]
sequence_ :: Monad m => [m a] -> m ()
\end{verbatim}}
\end{quote}
\item
A prefix ’m’ generalizes an existing function to a monadic form.
     Thus, for example:\par
\begin{quote}
{\haddockverb\begin{verbatim}
sum :: Num a => [a] -> a
msum :: MonadPlus m => [m a] -> m a
\end{verbatim}}
\end{quote}
\end{itemize}}
\subsection{Basic \haddockid{Monad} functions}
\begin{haddockdesc}
\item[\begin{tabular}{@{}l}
mapM :: (Traversable t, Monad m) => (a -> m b) -> t a -> m (t b)
\end{tabular}]
{\haddockbegindoc
Map each element of a structure to a monadic action, evaluate
 these actions from left to right, and collect the results. For
 a version that ignores the results see \haddockid{mapM{\char '137}}.\par
\subsubsection*{\textbf{Examples}}
\haddockid{mapM} is literally a \haddockid{traverse} with a type signature restricted
 to \haddockid{Monad}. Its implementation may be more efficient due to additional
 power of \haddockid{Monad}.\par}
\end{haddockdesc}
\begin{haddockdesc}
\item[\begin{tabular}{@{}l}
mapM{\char '137} :: (Foldable t, Monad m) => (a -> m b) -> t a -> m ()
\end{tabular}]
{\haddockbegindoc
Map each element of a structure to a monadic action, evaluate
 these actions from left to right, and ignore the results.  For a
 version that doesn't ignore the results see
 \haddockid{mapM}.\par
\haddockid{mapM{\char '137}} is just like \haddockid{traverse{\char '137}}, but specialised to monadic actions.\par}
\end{haddockdesc}
\begin{haddockdesc}
\item[\begin{tabular}{@{}l}
forM :: (Traversable t, Monad m) => t a -> (a -> m b) -> m (t b)
\end{tabular}]
{\haddockbegindoc
\haddockid{forM} is \haddockid{mapM} with its arguments flipped. For a version that
 ignores the results see \haddockid{forM{\char '137}}.\par}
\end{haddockdesc}
\begin{haddockdesc}
\item[\begin{tabular}{@{}l}
forM{\char '137} :: (Foldable t, Monad m) => t a -> (a -> m b) -> m ()
\end{tabular}]
{\haddockbegindoc
\haddockid{forM{\char '137}} is \haddockid{mapM{\char '137}} with its arguments flipped.  For a version that
 doesn't ignore the results see \haddockid{forM}.\par
\haddockid{forM{\char '137}} is just like \haddockid{for{\char '137}}, but specialised to monadic actions.\par}
\end{haddockdesc}
\begin{haddockdesc}
\item[\begin{tabular}{@{}l}
sequence :: (Traversable t, Monad m) => t (m a) -> m (t a)
\end{tabular}]
{\haddockbegindoc
Evaluate each monadic action in the structure from left to
 right, and collect the results. For a version that ignores the
 results see \haddockid{sequence{\char '137}}.\par
\subsubsection*{\textbf{Examples}}
Basic usage:\par
The first two examples are instances where the input and
 and output of \haddockid{sequence} are isomorphic.\par
\begin{quote}
{\haddockverb\begin{verbatim}
>>> sequence $ Right [1,2,3,4]
[Right 1,Right 2,Right 3,Right 4]

\end{verbatim}}
\end{quote}
\begin{quote}
{\haddockverb\begin{verbatim}
>>> sequence $ [Right 1,Right 2,Right 3,Right 4]
Right [1,2,3,4]

\end{verbatim}}
\end{quote}
The following examples demonstrate short circuit behavior
 for \haddockid{sequence}.\par
\begin{quote}
{\haddockverb\begin{verbatim}
>>> sequence $ Left [1,2,3,4]
Left [1,2,3,4]

\end{verbatim}}
\end{quote}
\begin{quote}
{\haddockverb\begin{verbatim}
>>> sequence $ [Left 0, Right 1,Right 2,Right 3,Right 4]
Left 0

\end{verbatim}}
\end{quote}}
\end{haddockdesc}
\begin{haddockdesc}
\item[\begin{tabular}{@{}l}
sequence{\char '137} :: (Foldable t, Monad m) => t (m a) -> m ()
\end{tabular}]
{\haddockbegindoc
Evaluate each monadic action in the structure from left to right,
 and ignore the results.  For a version that doesn't ignore the
 results see \haddockid{sequence}.\par
\haddockid{sequence{\char '137}} is just like \haddockid{sequenceA{\char '137}}, but specialised to monadic
 actions.\par}
\end{haddockdesc}
\begin{haddockdesc}
\item[\begin{tabular}{@{}l}
(=<<) :: Monad m => (a -> m b) -> m a -> m b
\end{tabular}]
{\haddockbegindoc
Same as \haddockid{>>=}, but with the arguments interchanged.\par
\begin{quote}
{\haddockverb\begin{verbatim}
as >>= f == f =<< as\end{verbatim}}
\end{quote}}
\end{haddockdesc}
\begin{haddockdesc}
\item[\begin{tabular}{@{}l}
(>=>) :: Monad m => (a -> m b) -> (b -> m c) -> a -> m c
\end{tabular}]
{\haddockbegindoc
Left-to-right composition of \haddockid{Kleisli} arrows.\par
'\haddocktt{(bs \haddockid{>=>} cs) a}' can be understood as the \haddocktt{do} expression\par
\begin{quote}
{\haddockverb\begin{verbatim}
do b <- bs a
   cs b
\end{verbatim}}
\end{quote}
or in terms of \haddocktt{\haddockid{(>>=)}} as\par
\begin{quote}
{\haddockverb\begin{verbatim}
bs a >>= cs\end{verbatim}}
\end{quote}}
\end{haddockdesc}
\begin{haddockdesc}
\item[\begin{tabular}{@{}l}
(<=<) :: Monad m => (b -> m c) -> (a -> m b) -> a -> m c
\end{tabular}]
{\haddockbegindoc
Right-to-left composition of Kleisli arrows. \haddocktt{(\haddockid{>=>})}, with the arguments
 flipped.\par
Note how this operator resembles function composition \haddocktt{(\haddockid{.})}:\par
\begin{quote}
{\haddockverb\begin{verbatim}
(.)   ::            (b ->   c) -> (a ->   b) -> a ->   c
(<=<) :: Monad m => (b -> m c) -> (a -> m b) -> a -> m c\end{verbatim}}
\end{quote}}
\end{haddockdesc}
\begin{haddockdesc}
\item[\begin{tabular}{@{}l}
forever :: Applicative f => f a -> f b
\end{tabular}]
{\haddockbegindoc
Repeat an action indefinitely.\par
\subsubsection*{\textbf{Examples}}
A common use of \haddockid{forever} is to process input from network sockets,
 \haddockid{Handle}s, and channels
 (e.g. \haddockid{MVar} and
 \haddockid{Chan}).\par
For example, here is how we might implement an \href{https://en.wikipedia.org/wiki/Echo_Protocol}{echo
 server}, using
 \haddockid{forever} both to listen for client connections on a network socket
 and to echo client input on client connection handles:\par
\begin{quote}
{\haddockverb\begin{verbatim}
echoServer :: Socket -> IO ()
echoServer socket = forever $ do
  client <- accept socket
  forkFinally (echo client) (\_ -> hClose client)
  where
    echo :: Handle -> IO ()
    echo client = forever $
      hGetLine client >>= hPutStrLn client
\end{verbatim}}
\end{quote}
Note that "forever" isn't necessarily non-terminating.
 If the action is in a \haddocktt{\haddockid{MonadPlus}} and short-circuits after some number of iterations.
 then \haddocktt{\haddockid{forever}} actually returns \haddockid{mzero}, effectively short-circuiting its caller.\par}
\end{haddockdesc}
\begin{haddockdesc}
\item[\begin{tabular}{@{}l}
void :: Functor f => f a -> f ()
\end{tabular}]
{\haddockbegindoc
\haddocktt{\haddockid{void} value} discards or ignores the result of evaluation, such
 as the return value of an \haddockid{IO} action.\par
\subsubsection*{\textbf{Examples}}
Replace the contents of a \haddocktt{\haddockid{Maybe} \haddockid{Int}} with unit:\par
\begin{quote}
{\haddockverb\begin{verbatim}
>>> void Nothing
Nothing

\end{verbatim}}
\end{quote}
\begin{quote}
{\haddockverb\begin{verbatim}
>>> void (Just 3)
Just ()

\end{verbatim}}
\end{quote}
Replace the contents of an \haddocktt{\haddockid{Either} \haddockid{Int} \haddockid{Int}}
 with unit, resulting in an \haddocktt{\haddockid{Either} \haddockid{Int} \haddocktt{()}}:\par
\begin{quote}
{\haddockverb\begin{verbatim}
>>> void (Left 8675309)
Left 8675309

\end{verbatim}}
\end{quote}
\begin{quote}
{\haddockverb\begin{verbatim}
>>> void (Right 8675309)
Right ()

\end{verbatim}}
\end{quote}
Replace every element of a list with unit:\par
\begin{quote}
{\haddockverb\begin{verbatim}
>>> void [1,2,3]
[(),(),()]

\end{verbatim}}
\end{quote}
Replace the second element of a pair with unit:\par
\begin{quote}
{\haddockverb\begin{verbatim}
>>> void (1,2)
(1,())

\end{verbatim}}
\end{quote}
Discard the result of an \haddockid{IO} action:\par
\begin{quote}
{\haddockverb\begin{verbatim}
>>> mapM print [1,2]
1
2
[(),()]

\end{verbatim}}
\end{quote}
\begin{quote}
{\haddockverb\begin{verbatim}
>>> void $ mapM print [1,2]
1
2

\end{verbatim}}
\end{quote}}
\end{haddockdesc}
\subsection{Generalisations of list functions}
\begin{haddockdesc}
\item[\begin{tabular}{@{}l}
join :: Monad m => m (m a) -> m a
\end{tabular}]
{\haddockbegindoc
The \haddockid{join} function is the conventional monad join operator. It
 is used to remove one level of monadic structure, projecting its
 bound argument into the outer level.\par
'\haddocktt{\haddockid{join} bss}' can be understood as the \haddocktt{do} expression\par
\begin{quote}
{\haddockverb\begin{verbatim}
do bs <- bss
   bs
\end{verbatim}}
\end{quote}
\subsubsection*{\textbf{Examples}}
\begin{quote}
{\haddockverb\begin{verbatim}
>>> join [[1, 2, 3], [4, 5, 6], [7, 8, 9]]
[1,2,3,4,5,6,7,8,9]

\end{verbatim}}
\end{quote}
\begin{quote}
{\haddockverb\begin{verbatim}
>>> join (Just (Just 3))
Just 3

\end{verbatim}}
\end{quote}
A common use of \haddockid{join} is to run an \haddockid{IO} computation returned from
 an \haddockid{STM} transaction, since \haddockid{STM} transactions
 can't perform \haddockid{IO} directly. Recall that\par
\begin{quote}
{\haddockverb\begin{verbatim}
atomically :: STM a -> IO a
\end{verbatim}}
\end{quote}
is used to run \haddockid{STM} transactions atomically. So, by
 specializing the types of \haddockid{atomically} and \haddockid{join} to\par
\begin{quote}
{\haddockverb\begin{verbatim}
atomically :: STM (IO b) -> IO (IO b)
join       :: IO (IO b)  -> IO b
\end{verbatim}}
\end{quote}
we can compose them as\par
\begin{quote}
{\haddockverb\begin{verbatim}
join . atomically :: STM (IO b) -> IO b
\end{verbatim}}
\end{quote}
to run an \haddockid{STM} transaction and the \haddockid{IO} action it
 returns.\par}
\end{haddockdesc}
\begin{haddockdesc}
\item[\begin{tabular}{@{}l}
msum :: (Foldable t, MonadPlus m) => t (m a) -> m a
\end{tabular}]
{\haddockbegindoc
The sum of a collection of actions using \haddockid{(<|>)}, generalizing \haddockid{concat}.\par
\haddockid{msum} is just like \haddockid{asum}, but specialised to \haddockid{MonadPlus}.\par
\subsubsection*{\textbf{Examples}}
Basic usage, using the \haddockid{MonadPlus} instance for \haddockid{Maybe}:\par
\begin{quote}
{\haddockverb\begin{verbatim}
>>> msum [Just "Hello", Nothing, Just "World"]
Just "Hello"

\end{verbatim}}
\end{quote}}
\end{haddockdesc}
\begin{haddockdesc}
\item[\begin{tabular}{@{}l}
filterM :: Applicative m => (a -> m Bool) -> {\char 91}a{\char 93} -> m {\char 91}a{\char 93}
\end{tabular}]
{\haddockbegindoc
This generalizes the list-based \haddockid{filter} function.\par
\begin{quote}
{\haddockverb\begin{verbatim}
runIdentity (filterM (Identity . p) xs) == filter p xs\end{verbatim}}
\end{quote}
\subsubsection*{\textbf{Examples}}
\begin{quote}
{\haddockverb\begin{verbatim}
>>> filterM (\x -> do
      putStrLn ("Keep: " ++ show x ++ "?")
      answer <- getLine
      pure (answer == "y"))
    [1, 2, 3]
Keep: 1?
y
Keep: 2?
n
Keep: 3?
y
[1,3]

\end{verbatim}}
\end{quote}
\begin{quote}
{\haddockverb\begin{verbatim}
>>> filterM (\x -> do
      putStr (show x)
      x' <- readLn
      pure (x == x'))
    [1, 2, 3]
12
22
33
[2,3]

\end{verbatim}}
\end{quote}}
\end{haddockdesc}
\begin{haddockdesc}
\item[\begin{tabular}{@{}l}
mapAndUnzipM :: Applicative m => (a -> m (b, c)) -> {\char 91}a{\char 93} -> m ({\char 91}b{\char 93}, {\char 91}c{\char 93})
\end{tabular}]
{\haddockbegindoc
The \haddockid{mapAndUnzipM} function maps its first argument over a list, returning
 the result as a pair of lists. This function is mainly used with complicated
 data structures or a state monad.\par}
\end{haddockdesc}
\begin{haddockdesc}
\item[\begin{tabular}{@{}l}
zipWithM :: Applicative m => (a -> b -> m c) -> {\char 91}a{\char 93} -> {\char 91}b{\char 93} -> m {\char 91}c{\char 93}
\end{tabular}]
{\haddockbegindoc
The \haddockid{zipWithM} function generalizes \haddockid{zipWith} to arbitrary applicative functors.\par}
\end{haddockdesc}
\begin{haddockdesc}
\item[\begin{tabular}{@{}l}
zipWithM{\char '137} :: Applicative m => (a -> b -> m c) -> {\char 91}a{\char 93} -> {\char 91}b{\char 93} -> m ()
\end{tabular}]
{\haddockbegindoc
\haddockid{zipWithM{\char '137}} is the extension of \haddockid{zipWithM} which ignores the final result.\par}
\end{haddockdesc}
\begin{haddockdesc}
\item[\begin{tabular}{@{}l}
foldM :: (Foldable t, Monad m) => (b -> a -> m b) -> b -> t a -> m b
\end{tabular}]
{\haddockbegindoc
The \haddockid{foldM} function is analogous to \haddockid{foldl}, except that its result is
encapsulated in a monad. Note that \haddockid{foldM} works from left-to-right over
the list arguments. This could be an issue where \haddocktt{(\haddockid{>>})} and the `folded
function' are not commutative.\par
\begin{quote}
{\haddockverb\begin{verbatim}
foldM f a1 [x1, x2, ..., xm]

==

do
  a2 <- f a1 x1
  a3 <- f a2 x2
  ...
  f am xm\end{verbatim}}
\end{quote}
If right-to-left evaluation is required, the input list should be reversed.\par
Note: \haddockid{foldM} is the same as \haddockid{foldlM}\par}
\end{haddockdesc}
\begin{haddockdesc}
\item[\begin{tabular}{@{}l}
foldM{\char '137} :: (Foldable t, Monad m) => (b -> a -> m b) -> b -> t a -> m ()
\end{tabular}]
{\haddockbegindoc
Like \haddockid{foldM}, but discards the result.\par}
\end{haddockdesc}
\begin{haddockdesc}
\item[\begin{tabular}{@{}l}
replicateM :: Applicative m => Int -> m a -> m {\char 91}a{\char 93}
\end{tabular}]
{\haddockbegindoc
\haddocktt{\haddockid{replicateM} n act} performs the action \haddocktt{act} \haddocktt{n} times,
 and then returns the list of results.\par
\begin{quote}
{\haddockverb\begin{verbatim}
replicateM n (pure x) == replicate n x\end{verbatim}}
\end{quote}
\subsubsection*{\textbf{Examples}}
\begin{quote}
{\haddockverb\begin{verbatim}
>>> replicateM 3 getLine
hi
heya
hiya
["hi","heya","hiya"]

\end{verbatim}}
\end{quote}
\begin{quote}
{\haddockverb\begin{verbatim}
>>> import Control.Monad.State

>>> runState (replicateM 3 $ state $ \s -> (s, s + 1)) 1
([1,2,3],4)

\end{verbatim}}
\end{quote}}
\end{haddockdesc}
\begin{haddockdesc}
\item[\begin{tabular}{@{}l}
replicateM{\char '137} :: Applicative m => Int -> m a -> m ()
\end{tabular}]
{\haddockbegindoc
Like \haddockid{replicateM}, but discards the result.\par
\subsubsection*{\textbf{Examples}}
\begin{quote}
{\haddockverb\begin{verbatim}
>>> replicateM_ 3 (putStr "a")
aaa

\end{verbatim}}
\end{quote}}
\end{haddockdesc}
\subsection{Conditional execution of monadic expressions}
\begin{haddockdesc}
\item[\begin{tabular}{@{}l}
guard :: Alternative f => Bool -> f ()
\end{tabular}]
{\haddockbegindoc
Conditional failure of \haddockid{Alternative} computations. Defined by\par
\begin{quote}
{\haddockverb\begin{verbatim}
guard True  = pure ()
guard False = empty
\end{verbatim}}
\end{quote}
\subsubsection*{\textbf{Examples}}
Common uses of \haddockid{guard} include conditionally signalling an error in
 an error monad and conditionally rejecting the current choice in an
 \haddockid{Alternative}-based parser.\par
As an example of signalling an error in the error monad \haddockid{Maybe},
 consider a safe division function \haddocktt{safeDiv x y} that returns
 \haddockid{Nothing} when the denominator \haddocktt{y} is zero and \haddocktt{\haddockid{Just} (x `div`\\ y)} otherwise. For example:\par
\begin{quote}
{\haddockverb\begin{verbatim}
>>> safeDiv 4 0
Nothing

\end{verbatim}}
\end{quote}
\begin{quote}
{\haddockverb\begin{verbatim}
>>> safeDiv 4 2
Just 2

\end{verbatim}}
\end{quote}
A definition of \haddocktt{safeDiv} using guards, but not \haddockid{guard}:\par
\begin{quote}
{\haddockverb\begin{verbatim}
safeDiv :: Int -> Int -> Maybe Int
safeDiv x y | y /= 0    = Just (x `div` y)
            | otherwise = Nothing
\end{verbatim}}
\end{quote}
A definition of \haddocktt{safeDiv} using \haddockid{guard} and \haddockid{Monad} \haddocktt{do}-notation:\par
\begin{quote}
{\haddockverb\begin{verbatim}
safeDiv :: Int -> Int -> Maybe Int
safeDiv x y = do
  guard (y /= 0)
  return (x `div` y)
\end{verbatim}}
\end{quote}}
\end{haddockdesc}
\begin{haddockdesc}
\item[\begin{tabular}{@{}l}
when :: Applicative f => Bool -> f () -> f ()
\end{tabular}]
{\haddockbegindoc
Conditional execution of \haddockid{Applicative} expressions. For example,\par
\subsubsection*{\textbf{Examples}}
\begin{quote}
{\haddockverb\begin{verbatim}
when debug (putStrLn "Debugging")\end{verbatim}}
\end{quote}
will output the string \haddocktt{Debugging} if the Boolean value \haddocktt{debug}
 is \haddockid{True}, and otherwise do nothing.\par
\begin{quote}
{\haddockverb\begin{verbatim}
>>> putStr "pi:" >> when False (print 3.14159)
pi:

\end{verbatim}}
\end{quote}}
\end{haddockdesc}
\begin{haddockdesc}
\item[\begin{tabular}{@{}l}
unless :: Applicative f => Bool -> f () -> f ()
\end{tabular}]
{\haddockbegindoc
The reverse of \haddockid{when}.\par
\subsubsection*{\textbf{Examples}}
\begin{quote}
{\haddockverb\begin{verbatim}
>>> do x <- getLine
       unless (x == "hi") (putStrLn "hi!")
comingupwithexamplesisdifficult
hi!

\end{verbatim}}
\end{quote}
\begin{quote}
{\haddockverb\begin{verbatim}
>>> unless (pi > exp 1) Nothing
Just ()

\end{verbatim}}
\end{quote}}
\end{haddockdesc}
\subsection{Monadic lifting operators}
\begin{haddockdesc}
\item[\begin{tabular}{@{}l}
liftM :: Monad m => (a1 -> r) -> m a1 -> m r
\end{tabular}]
{\haddockbegindoc
Promote a function to a monad.
 This is equivalent to \haddockid{fmap} but specialised to Monads.\par}
\end{haddockdesc}
\begin{haddockdesc}
\item[\begin{tabular}{@{}l}
liftM2 :: Monad m => (a1 -> a2 -> r) -> m a1 -> m a2 -> m r
\end{tabular}]
{\haddockbegindoc
Promote a function to a monad, scanning the monadic arguments from
 left to right.\par
\subsubsection*{\textbf{Examples}}
\begin{quote}
{\haddockverb\begin{verbatim}
>>> liftM2 (+) [0,1] [0,2]
[0,2,1,3]

\end{verbatim}}
\end{quote}
\begin{quote}
{\haddockverb\begin{verbatim}
>>> liftM2 (+) (Just 1) Nothing
Nothing

\end{verbatim}}
\end{quote}
\begin{quote}
{\haddockverb\begin{verbatim}
>>> liftM2 (+) (+ 3) (* 2) 5
18

\end{verbatim}}
\end{quote}}
\end{haddockdesc}
\begin{haddockdesc}
\item[\begin{tabular}{@{}l}
liftM3 :: Monad m => (a1 -> a2 -> a3 -> r) -> m a1 -> m a2 -> m a3 -> m r
\end{tabular}]
{\haddockbegindoc
Promote a function to a monad, scanning the monadic arguments from
 left to right (cf. \haddockid{liftM2}).\par}
\end{haddockdesc}
\begin{haddockdesc}
\item[\begin{tabular}{@{}l}
liftM4 :: Monad m => (a1 -> a2 -> a3 -> a4 -> r) -> m a1 -> m a2 -> m a3 -> m a4 -> m r
\end{tabular}]
{\haddockbegindoc
Promote a function to a monad, scanning the monadic arguments from
 left to right (cf. \haddockid{liftM2}).\par}
\end{haddockdesc}
\begin{haddockdesc}
\item[\begin{tabular}{@{}l}
liftM5 :: Monad m =>\\\ \ \ \ \ \ \ \ \ \ (a1 -> a2 -> a3 -> a4 -> a5 -> r) -> m a1 -> m a2 -> m a3 -> m a4 -> m a5 -> m r
\end{tabular}]
{\haddockbegindoc
Promote a function to a monad, scanning the monadic arguments from
 left to right (cf. \haddockid{liftM2}).\par}
\end{haddockdesc}
\begin{haddockdesc}
\item[\begin{tabular}{@{}l}
ap :: Monad m => m (a -> b) -> m a -> m b
\end{tabular}]
{\haddockbegindoc
In many situations, the \haddockid{liftM} operations can be replaced by uses of
\haddockid{ap}, which promotes function application.\par
\begin{quote}
{\haddockverb\begin{verbatim}
return f `ap` x1 `ap` ... `ap` xn\end{verbatim}}
\end{quote}
is equivalent to\par
\begin{quote}
{\haddockverb\begin{verbatim}
liftM<n> f x1 x2 ... xn\end{verbatim}}
\end{quote}
\subsubsection*{\textbf{Examples}}
\begin{quote}
{\haddockverb\begin{verbatim}
>>> pure (\x y z -> x + y * z) `ap` Just 1 `ap` Just 5 `ap` Just 10
Just 51

\end{verbatim}}
\end{quote}}
\end{haddockdesc}